% NbTi-to-NbTi superconducting joint: conceptual cross-section + build steps.
\documentclass{article}
\usepackage[paperwidth=120cm,paperheight=120cm,margin=10mm]{geometry}
\pagestyle{empty}
\usepackage{fontspec}
\usepackage[bidi=basic]{babel}
\babelprovide[import]{english}
\babelprovide[main, import]{persian}
\babelfont{rm}[Renderer=Node, Path=../fonts/, Extension=.ttf, UprightFont=*-Regular, BoldFont=*-Bold]{Vazirmatn}
\babelfont{sf}[Renderer=Node, Path=../fonts/, Extension=.ttf, UprightFont=*-Regular, BoldFont=*-Bold]{Vazirmatn}
\providecommand{\setRTL}{}\providecommand{\setLTR}{}
\providecommand{\RL}[1]{\foreignlanguage{persian}{#1}}
\providecommand{\LR}[1]{\babelsublr{#1}}
\usepackage{tikz}
\usepackage{pgfplots}
\pgfplotsset{compat=1.18}
\usetikzlibrary{arrows.meta, positioning, shapes.geometric, shapes.misc, calc, fit, backgrounds, decorations.pathmorphing, patterns, angles, quotes}
\begin{document}
\setRTL
\tikzset{every node/.append style={execute at begin node=\setRTL}}
\babelsublr{\begin{tikzpicture}[>=Stealth, font=\small]

  \node[font=\bfseries] at (3.4,3.1) {اتصال NbTi--NbTi (مفهومی)};

  % ----- left: conceptual cross-section of the joint -----
  % copper stabilizer of wire 1 (etched away at the joint region)
  \fill[orange!30] (-0.2,1.4) rectangle (2.0,2.0);
  \fill[orange!30] (-0.2,0.4) rectangle (2.0,1.0);
  \node[font=\footnotesize, orange!60!black] at (0.9,2.25) {پایدارساز Cu (سیم ۱)};
  \node[font=\footnotesize, orange!60!black] at (0.9,0.15) {پایدارساز Cu (سیم ۲)};

  % exposed bare NbTi filaments brought side by side, inside solder
  \fill[gray!25, rounded corners=3pt] (2.0,0.55) rectangle (5.4,1.85);
  \node[font=\footnotesize, gray!50!black] at (4.4,1.65) {لحیم ابررسانا (PbBi/PbSn)};
  % filaments
  \foreach \y in {1.05,1.2,1.35}{
    \draw[very thick, blue!60!black] (2.05,\y) -- (5.0,\y);
  }
  \foreach \y in {0.85,1.0}{
    \draw[very thick, red!65!black] (2.05,\y) -- (5.0,\y);
  }
  \node[font=\footnotesize, blue!60!black] at (3.5,1.5) {رشته‌های لخت NbTi};

  % joint box outline
  \draw[thick, black!70, dashed, rounded corners=4pt] (1.9,0.35) rectangle (5.6,2.05);
  \node[font=\footnotesize, black!70] at (5.05,0.2) {جعبهٔ اتصال};

  % ----- right: ordered build steps -----
  \def\bx{6.6}
  \def\dy{0.72}
  \node[font=\bfseries, anchor=west] at (\bx,2.85) {مراحل ساخت};
  \node[anchor=west, text width=6.0cm] at (\bx,2.85-1*\dy) {۱) اچ اسیدی مس برای نمایان‌کردن NbTi لخت};
  \node[anchor=west, text width=6.0cm] at (\bx,2.85-2*\dy) {۲) قرار دادن رشته‌های هر دو سیم کنار هم};
  \node[anchor=west, text width=6.0cm] at (\bx,2.85-3*\dy) {۳) لحیم با آلیاژ ابررسانا (ابررسانا در $4.2\,$K)};
  \node[anchor=west, text width=6.0cm] at (\bx,2.85-4*\dy) {\quad\itshape یا ۴) جوش سرد / نقطه‌ای مستقیم NbTi};
  \node[anchor=west, text width=6.0cm] at (\bx,2.85-5*\dy) {۵) محصورسازی در جعبهٔ اتصال (پوشش مکانیکی و مغناطیسی)};

\end{tikzpicture}}
\end{document}
